\documentclass[a4paper,11pt]{article}
\usepackage{plantilla}

\titulo{Tortilla de patatas}
\autor{Carlos Sesma}

% Una plantilla puede añadir lo suyo encima de plantilla.sty: aquí, una ficha
% con los datos de cabecera de la receta.
\newcommand{\ficha}[3]{%
    \begin{center}
        \begin{tabular}{l@{\quad}l@{\quad}l}
            \textbf{Raciones} & \textbf{Preparación} & \textbf{Cocción} \\
            #1                & #2                   & #3               \\
        \end{tabular}
    \end{center}
}

\begin{document}

\begin{center}
    {\LARGE\bfseries \eltitulo}
\end{center}

\ficha{4}{20 min}{25 min}

\section{Ingredientes}

\begin{itemize}
    \item 6 patatas medianas
    \item 6 huevos
    \item 1 cebolla (opcional, según bando)
    \item Aceite de oliva y sal
\end{itemize}

\section{Elaboración}

\begin{enumerate}
    \item Pelar y cortar las patatas en láminas finas. Salar.
    \item Confitarlas a fuego suave en abundante aceite, sin que lleguen a
          dorarse, unos 20 minutos.
    \item Escurrir bien y mezclar con los huevos batidos. Dejar reposar
          5 minutos.
    \item Cuajar en la sartén a fuego medio, dando la vuelta una vez.
\end{enumerate}

\begin{nota}[Truco]
    El reposo antes de cuajar es lo que une la patata con el huevo. Saltárselo
    es la causa habitual de que la tortilla se desmonte al girarla.
\end{nota}

\begin{aviso}[Cuidado]
    Si el aceite humea, la patata se fríe en vez de confitarse y la textura
    cambia por completo.
\end{aviso}

\end{document}
